\section{Introduction}
\label{sec:intro}

Systems that serve many user queries often select one large language model (LLM) from a pool to balance \textbf{response quality} against \textbf{inference cost}.
Recent \textbf{multi-LLM routing} systems---RouteLLM, HybridLLM, GraphRouter, RouterBench, and others---typically \textbf{supervise} routing: they train routers on human preferences, historical win/loss outcomes, reward-model scores, or task labels to predict which model should answer each query~\citep{ong2024routellm,ding2024hybrid,feng2025graphrouter,hu2024routerbench}.
These approaches have shown that adaptive selection can outperform always using a single model, especially when pool members exhibit \textbf{complementary strengths} on different inputs~\citep{hu2024routerbench,chen2024frugalgpt}.

Yet a prior question remains largely unanswered:
\emph{what routing-relevant information exists in signals computed \textbf{without} routing supervision at extraction time?}
By \textbf{unsupervised signals} we mean features derived from the query and from model responses using fixed, pre-specified statistics---not representations learned end-to-end from routing datasets (Section~\ref{sec:problem} defines this precisely).
Labels may still be used on a \textbf{calibration} split to test whether those signals predict \textbf{routing need} and to fit combination weights; that supervision applies to \emph{analysis} and \emph{policy fit}, not to inventing the signals themselves.

\paragraph{Motivation.}
We consider a fixed \textbf{model pool} whose members have \textbf{distinct capabilities} (not interchangeable).
For the primary routing study we use a designated lower-capability member $\Mlo$ and higher-capability member $\Mhi$.
For many queries both succeed; for others $\Mlo$ fails while $\Mhi$ would succeed---these are \textbf{opportunity} queries where escalation is warranted.
Prior routing work assumes that a supervised router will learn to identify such cases; we instead ask whether \textbf{query-derived}, \textbf{model-response}, and \textbf{cross-model comparative} signals---entropy, confidence, disagreement, and related quantities~\citep{kadavath2022know,plaut2024calibration,he2023calibration,kuhn2023semantic}---already carry information about routing need before any routing classifier is trained.

\paragraph{Research gaps.}
Two gaps motivate our study.
\textbf{Gap 1 --- Supervised routing dominates.}
Most LLM routers \emph{define} routing through labeled data: preferences, outcomes, or learned query encoders trained on quality scores~\citep{ong2024routellm,ding2024hybrid,feng2025graphrouter,zhang2026agentrouter}.
It remains unclear which \textbf{unsupervised} signal types predict when a lower-capability pool member fails but a higher-capability member would succeed, and whether such signals can \textbf{enable routing decisions} without routing-specific training at extraction time.
\textbf{Gap 2 --- Signal types are not tested systematically.}
Uncertainty and calibration work studies whether models ``know what they know''~\citep{kadavath2022know,plaut2024calibration}, but prior routing papers rarely perform \textbf{nested ablations} over query-derived, model-response, and cross-model comparative signals against a single routing-need target.
Understanding \emph{which} signal layer carries information---and \emph{whether} that information translates into better accuracy--cost routing---requires separating \textbf{signal analysis} from \textbf{routing evaluation}.

\paragraph{Approach.}
We study \emph{routing from unsupervised signals} in a fixed capability-differentiated pool.
Our pipeline (Figure~\ref{fig:pipeline}; Section~\ref{sec:method}) deliberately separates two scientific threads:
(1)~\textbf{signal understanding}---which unsupervised signals predict routing need $r(q)$ on calib (hypotheses H1--H3); and
(2)~\textbf{routing performance}---whether a threshold policy built from frozen features improves accuracy--cost trade-offs on test (hypothesis H4).
We extract candidate signals without $r(q)$, analyze nested signal types against offline routing need, freeze a compact feature vector $x(q)$, fit weights $(\lambda,\tau)$ on calib only, and evaluate with \textbf{Pareto} curves over task accuracy and average inference cost.
Negative results on any hypothesis are valid findings: informative signals may coexist with poor routing, or vice versa.

\paragraph{Contributions.}
\begin{enumerate}
  \item \textbf{Formulation.}
    Routing need $r(q)$, signal assumptions and taxonomy, and a precise \textbf{supervised vs.\ unsupervised} contrast at signal extraction (Sections~\ref{sec:problem}--\ref{sec:signal-extraction}).
  \item \textbf{Signal analysis.}
    Nested statistical tests (H1--H3) of whether query-derived, model-response, and cross-model comparative signals predict $r(q)$ on calib (Section~\ref{sec:results-signals}).
  \item \textbf{Signal selection.}
    A frozen $x(q)$ chosen from calib analysis before any routing claims (Section~\ref{sec:signal-selection}).
  \item \textbf{Routing policies.}
    Hand-designed and learned combinations of signals via $(\lambda,\tau)$ (Section~\ref{sec:routing-policy}).
  \item \textbf{Pareto evaluation.}
    Accuracy and cost reported jointly on test---not a single opaque ``better router'' score (Section~\ref{sec:results-routing}).
\end{enumerate}

We do not reproduce end-to-end supervised routers; we characterize what unsupervised signals contribute and test simple policies built from them.
Our main contrast is \textbf{supervised vs.\ unsupervised at signal extraction}, not agent orchestration or pre- vs post-inference timing as the headline claim.
